% Generated draft: requires mathematical and source review.
\documentclass[11pt]{article}
\usepackage[margin=25mm]{geometry}
\usepackage{fontspec}
\setmainfont{Latin Modern Roman}
\usepackage{amsmath,amssymb,mathtools,unicode-math}
\setmathfont{Latin Modern Math}
\usepackage{enumitem,needspace,xurl,hyperref}
\setlength{\parindent}{0pt}
\setlength{\parskip}{5pt}
\setlength{\emergencystretch}{3em}
\newcommand{\N}{\mathbb{N}}
\newcommand{\Z}{\mathbb{Z}}
\newcommand{\Q}{\mathbb{Q}}
\newcommand{\R}{\mathbb{R}}
\newcommand{\C}{\mathbb{C}}
\newcommand{\F}{\mathbb{F}}
\newcommand{\A}{\mathbb{A}}
\newcommand{\PP}{\mathbb{P}}
\newcommand{\E}{\mathbb{E}}
\newcommand{\eps}{\varepsilon}
\newcommand{\dd}{\,\mathrm{d}}
\begin{document}
\section*{UnsolvedMath ID: 30004996}
\textbf{Source problem code:} catalogue entry
\subsection{Definitions and mathematical statement}
The supplied record does not specify the allowed group or alphabet parameters, or define hyper-aperiodicity. The following is therefore a conditional formulation with explicit conventions, requiring comparison with the source.

Let \(\Gamma\) be a countable discrete group with identity \(e\), and let \(k\) and \(q\) be finite nonempty discrete alphabets. These parameter conventions are not explicit in the supplied statement; in particular, any intended infinitude assumption on \(\Gamma\) or lower bound on \(|k|\) remains unresolved. Equip \(k^\Gamma\) and \(q^\Gamma\) with their product topologies and the left shift action
\[
(g\cdot x)(h)=x(g^{-1}h)\qquad(g,h\in\Gamma).
\]
For either alphabet \(a\), define
\[
\operatorname{Free}(a^\Gamma)=\{x\in a^\Gamma: g\cdot x\ne x\text{ for every }g\in\Gamma\setminus\{e\}\}.
\]
Write \(X=\operatorname{Free}(k^\Gamma)\), with its inherited Borel structure. A subshift of finite type \(S\subseteq q^\Gamma\) is a set for which there are a finite set \(D\subseteq\Gamma\) and a set of allowed patterns \(P\subseteq q^D\) such that
\[
S=\{y\in q^\Gamma:(g\cdot y)|_D\in P\text{ for every }g\in\Gamma\}.
\]
For a Borel map \(f:X\to q\), define its equivariant coding map by
\[
\pi_f(x)(g)=f(g^{-1}\cdot x).
\]
An \(S\)-coloring means that \(\pi_f(X)\subseteq S\); it is aperiodic when
\[
\pi_f(X)\subseteq\operatorname{Free}(q^\Gamma).
\]

Under the interpretation that a hyper-aperiodic coloring has a coding image whose entire closure consists of free points, the question is whether, for every choice of the intended parameters and every such \(S\),
\[
\begin{gathered}
\exists\text{ Borel }f:X\to q\quad
\pi_f(X)\subseteq S\cap\operatorname{Free}(q^\Gamma)
\\[2pt]
\Longrightarrow
\\[2pt]
\exists\text{ Borel }b:X\to q\quad
\overline{\pi_b(X)}\subseteq S\cap\operatorname{Free}(q^\Gamma),
\end{gathered}
\]
where closure is taken in \(q^\Gamma\). The desired coloring \(b\) need not equal the given coloring \(f\).

The source terminology must be checked: if hyper-aperiodicity instead means only that each individual coded configuration has a free orbit closure, the conclusion would be
\[
\pi_b(X)\subseteq S,
\qquad
\overline{\Gamma\cdot\pi_b(x)}\subseteq\operatorname{Free}(q^\Gamma)
\quad\text{for every }x\in X.
\]
This pointwise condition must not be silently identified with the condition on the closure of the whole coding image.

The supplied contextual discussion says that merely replacing continuity by Borel measurability in its referenced Theorem 3 is insufficient and that further assumptions are necessary. It presents this concrete question as a test case related to its Problem 1; neither that theorem nor that broader problem is supplied here.
\subsection{Short English statement}
Does the existence of a Borel aperiodic coloring satisfying finite local constraints imply the existence of a Borel hyper-aperiodic coloring satisfying the same constraints on \(\operatorname{Free}(k^\Gamma)\)? The group and alphabet ranges, and the precise closure condition intended by hyper-aperiodicity, require source confirmation. This is a statement-only draft requiring human mathematical review.
\subsection{Sources}
\begin{itemize}
\item Local UnsolvedMath registry, permanent ID \(30004996\), source code \texttt{catalogue entry}, title \emph{Hyper-Aperiodic Colorings from Aperiodic SFT Colorings}. The supplied original excerpt labels the question Problem 2. No individual proposer is identified; the recorded proposed year is 2022.
\item Supplied source citation: \emph{Groups and Dynamics: Topology, Measure, and Borel Structure} (2023), Oberwolfach Reports, pp.\ 107--169, DOI \texttt{10.4171/owr/2022/3}. \url{https://doi.org/10.4171/owr/2022/3}. The literature-source list also labels this reference \emph{Oberwolfach Report 3/2022}; the differing year labels are retained without resolution.
\item Supplied literature source: Bernshteyn, \emph{Equivariant maps to subshifts whose points have small stabilizers}. \url{https://arxiv.org/abs/2106.09673}.
\item Literature assessment dated exactly \texttt{2026-08-22}: the supplied curation labels the question open and says that existing theorems require continuity or extra structural assumptions. This is supplied dated curation, not an independently verified literature assessment. The background explicitly requests independent verification of time-sensitive status and cited claims.
\item The registry labels the statement \texttt{corrected\_verified} and reports reconstruction and checking against the official Oberwolfach report. These are supplied provenance claims; no independent source consultation or mathematical verification is asserted here.
\end{itemize}
\end{document}
